\section{Compact companion: a squared transport coordinate}
\label{sec:compact-companion}

The direct-calculus constructions and the following compact companion
optimize different quantities.  This section gives two literal polynomials
in ten \emph{strictly positive} integer variables, each of exact degree
$232$ and $85$ binary arithmetic operations.  The trace parent has
$87$ operations at degree $232$; the inverse-Lucas parent has $85$
operations at degree $304$.  The new squared-coordinate variant also
reduces the degree of one carrier numerator from $57$ to $21$.
The total degree remains $232$.

\paragraph{Late comparison and attribution.}
At the later read-only snapshot
\texttt{0d75e561a8cefc8c30fc4224605a8984c810ad1c}, the incoming
Pooled Powers Separated Radicals archive already contains exactly the
same $85$-gate $bV$ polynomial: all rows, all ten inputs and all four fixed
coefficient values agree.  The $85/232$ metric point is therefore not
claimed as new here.  That equation is a reproduced baseline.  The
distinct contribution is the single replacement $L=bV\mapsto b^2$
and its soundness/completeness proof, which lowers the carrier degree
without changing total degree or operation count.  The exact comparison
is recorded in the companion's \texttt{LATE\_COMPARISON.json}.

The two parents are packages \texttt{163-trace-radical} and
\texttt{164-unsquared-factors-linear-selectors} of the pinned ProveIt
snapshot.  The former explicitly left open the squared-coordinate
transport argument.  Its difficulty was a correction growing with a
negative rank.  Evaluating the inverse correction at a positive Pell
coordinate removes that difficulty.  This proves the particular new
variant below; it does not assert worldwide novelty or global minimality.
This compact component retains the native/K215 ancestry of its fixed
coefficient tuple.

\subsection{Complete polynomial}

The quantified inputs are
\[
 J,F,\alpha,i,T,w,s,n,\eta,\zeta\in\mathbb Z_{>0}.
\]
The four fixed integers $\beta,K_0,M_C,M_F$ are the complete indexed
BFS943/radix67 tuple, reconstructed from the supplied finite recipe.
In particular $\beta=2^{590069}-1$; the largest port has $592884$ bits.
All four actual ports, rather than arbitrary integers with the same
elementary bounds, are needed for the logical calibration.
Define polynomial abbreviations
\begin{align*}
 h&=\beta J,&q&=h+1,& C&=q-F-\alpha,\\
 X&=qw,&Y&=qs,&k&=\eta+\zeta,\\
 c&=kY+\eta,&A&=Y(X+1)+1,&\Delta&=A^2-1,\\
 H&=4A-5,&v&=XY^2,&\ell&=vk,\\
 E&=1+\ell(\ell+k),&&&
\end{align*}
and
\begin{align*}
 R&=[h(q-F)+\alpha](q^2-1)+(M_C+qM_F)J,\\
 B&=(K_0+w)C-F,\\
 U&=ic^2,\qquad a=\Delta U^2E,\qquad F_0=1+a,\qquad G=EF_0,\\
 f&=1+2a,\qquad Q=4\Delta aF_0,\\
 b&=cT-R,\qquad V=fb-c,\qquad K=X(3c-2X)+2,\\
 e_6&=6(k-R-1)-\ell(k^2-4),\qquad D_6=6q^6,\\
 M&=hD_6H,\\
 Z&=hD_6[H(n-LB)+K]-H[D_6B+h(e_6-6q^4s)] .
\end{align*}
The principal weight is $L=b^2$; a companion uses $L=bV$.
The equation is
\begin{equation}
 \boxed{\;Q\bigl(Z^2-M^2GV^2\bigr)+\bigl(M^2G-Z^2\bigr)=0.\;}
 \label{eq:compact232}
\end{equation}
Every symbol on the right is a polynomial abbreviation.  There are no
additional quantified variables, variable exponents, recurrence operators,
or uncomputed coefficients.  The literal circuit computes
$\ell=(XY)(kY)$ and $f=a+F_0$, thereby sharing existing operations.

Composed with the named inherited native/refutation calibration, the
ordinary mathematical theorem is
\[
 \mathrm{ZFC}\vdash
 \left(\neg\exists\boldsymbol{x}\in\mathbb Z_{>0}^{10}\;
 P(\boldsymbol{x})=0\right)
 \longleftrightarrow\operatorname{Con}(\mathrm{ZFC}).
\]
The domain is essential.  If nonnegative variables are desired, substitute
$x_j=u_j+1$ for each of the ten inputs.  This leaves the dimension and
total degree unchanged, though the displayed operation count refers to
the positive-input circuit.

\subsection{The new transport proof}

For the proof only, put
\[
 e_*={e_6\over6},\quad \rho={e_*\over q^6},\quad
 \sigma={s\over q^2},\quad t={B\over h},\quad
 \kappa={K\over H},\quad W=1+hL,\quad z={Z\over M}.
\]
No integrality of these quotients is initially assumed.  The exact
rational identity is
\begin{equation}
 z=n+\kappa+\sigma-Wt-\rho.\label{eq:compact-carrier}
\end{equation}

The trace parent's root extraction and low-index exclusions depend on
$Q,V,G$ and the final norm, not on $L$ or $e_6$.
At a positive zero they therefore still give
\[
 G=g^2,quad z=\varepsilon gy,quad
 QV^2-(Q-1)y^2=1,qquad g,y>0,quad\varepsilon\in\{-1,1\},
\]
with $z,y$ integral, and
\[
 b\ne0,\quad |b|>8f^3-\tfrac12,\quad
 g<f/2,\quad y<3|V|/2,\quad c<f/2.
\]
These statements are established before any sign assumption on $R$ or $B$.
In particular, since $b$ is a nonzero integer,
\[
 |V|\le f|b|+c<\tfrac32 f|b|,
 \qquad
 gy<\tfrac98 f^2|b|<\tfrac12 b^2.
\]
The last strict inequality follows from
$|b|>8f^3-1/2>(9/4)f^2$ for $f\ge1$.
Consequently $L=b^2$ satisfies
\[
 L>0,\qquad gy<L/2,\qquad |b|/q^6<L/2.
\]

The inverse correction has the decisive sign:
\[
 {vk(k^2-4)\over6}\ge0,
 \qquad e_*\le k-R-1<b,
\]
because $k\ge2$ and $b-(k-R-1)=cT-k+1>0$.
This holds for every integer $R$, including negative ranks.  Thus
$\rho<|b|/q^6$.  If $B\le-1$, then $Wt<-L$; using the positive
$n,\kappa,\sigma$ in \eqref{eq:compact-carrier} gives
\[
 n=\varepsilon gy-\kappa-\sigma+Wt+\rho
   <L/2-L+|b|/q^6<0,
\]
a contradiction.  Therefore $B\ge0$ and then $C>0$.
This is the new unbounded sign argument.  It is not inferred from
bounded testing or from a formal substitution into the forward-correction
proof.

\subsection{The remaining arithmetic order}

Once $C>0$, the inherited mask bound gives $3q<R<q^4-2$.
The unsquared trace obeys $f\equiv1\pmod c$, so the two auxiliary
residues remain
$V\equiv-c\pmod f$ and $V\equiv-R\pmod c$.
The retained signed-index theorem authenticates the main Pell rank as $R$.

For clarity, the inverse denominator argument must retain its order.
The first Pell recurrence has
\[
 k=2m+8vC_m+32v^2Z_m,qquad
 C_m=\binom{m+1}{3},\quad Z_m\in\mathbb N.
\]
Putting $d=2m-R-1$, $A_m=C_m+4vZ_m$ and
$S_m=k^2+2mk+4m^2-4$ gives the exact identity
\[
 3(e_*-d)=4v^2(24Z_m-A_mS_m).
\]
The carrier gives $q^6\mid(e_*-q^4s)$, hence $q^4\mid e_*$.
Since $v=q^3ws^2$, prime valuations in the displayed identity give
$q^4\mid(e_*-d)$ even at the prime $3$.
The bounds $|d|<R<q^4$ then force $d=0$.

It would be premature to cancel the remaining factor $3$ and infer
$q^6\mid e_*$.  Instead, the carrier first gives $H\mid6hK$.
The retained uniform ratio estimate and whole-product projection force
$X=2^R$.  Since $q\mid X$, the radix $q$ is a power of two.  Only now
does the identity imply $q^6\mid e_*$.  The carrier then gives
$q^2\mid s$ and $h\mid B$, supplying the exact native outer endpoint.
The parent's indexed native/refutation calibration completes soundness.

Conversely, a finite refutation gives the canonical outer witnesses and
the trace parent's positive auxiliary Pell pair.  Its actual positive
integer input is $T=(V+c+Rf)/(cf)$.
The inverse error can be negative, but the inherited explicit bounds give
\[
 |\rho|<vk^3/3<g/3,\qquad \kappa+\sigma<g/2.
\]
For either positive weight choose
\[
 n=gy+(1+hL)t+\rho-\kappa-\sigma.
\]
This is an integer and is strictly positive because $y\ge1$, $t\ge0$
and $g>\kappa+\sigma+|\rho|$.  The carrier now has $z=gy$, proving
completeness for all ten positive inputs.

\subsection{Exact complexity and scope}

Let $q_0=\beta J$, $k_0=\eta+\zeta$ and
$B_0=k_0sT-q_0^2(q_0-F)$.  The complete highest homogeneous component is
\[
 [P]_{232}=-9\,2^{10}q_0^{100}w^{26}s^{58}i^{10}k_0^{32}B_0^2.
\]
It is nonzero at the actual $\beta>0$.
For $L=b^2$, the degrees of $(Q,M,G,V,Z)$ are
$(80,11,50,40,21)$.  The term $QZ^2$ has degree $122$, while
$-QM^2GV^2$ has degree $232$ and remains the unique dominant summand.
For $L=bV$, $\deg Z=57$ and $\deg(QZ^2)=194$.

The reproducible finite checker covers every live gate by $36$ exact
free-polynomial cuts per graph, checks the complete $202$-monomial highest
component and all four coefficient payloads, and rejects $22$ deliberately
corrupted operators per graph.  Normal and optimized Python runs agree.
The finite recipe reconstructs every fixed integer and matches its pinned
raw SHA-256 identity.  These are algebraic artifact checks, separate from
the unbounded arguments and inherited logical calibration.
No new end-to-end proof-assistant development is claimed.
